\par\medskip
\hypertarget{native-074-article056}{}\pdfbookmark[2]{56.}{article-056}
\begin{center}56.\end{center}
Es lohnt sich, die vorangehenden Regeln an einem Beispiel zu erläutern. Es sei $\AscendingNode=172^0\,28^{\prime}\,13^{\prime\prime}{,}7$ und $i=34^0\,38^{\prime}\,1^{\prime\prime}{,}1$. Ferner sei die neue Ebene zum Äquator parallel, sodass $n=180^0$ gilt. Für den Winkel $\epsilon$, also die Schiefe der Ekliptik, setzen wir $=23^0\,27^{\prime}\,55^{\prime\prime}{,}8$. Damit erhalten wir
\begin{center}
\begin{tabular}{@{}ll@{}}
$\AscendingNode-n=$ & $-7^0\,31^{\prime}\,46^{\prime\prime}{,}3$\\[.3em]
$i+\epsilon=$ & $58\quad5\quad56{,}9$\\[.3em]
$i-\epsilon=$ & $11\quad10\quad5{,}3$\\[.3em]
\makebox[18em][l]{$\log\sin\tfrac{1}{2}(\AscendingNode-n)$\hfill .\,.} & $8{,}817\,3026_n$\\[.3em]
\makebox[18em][l]{$\log\sin\tfrac{1}{2}(i+\epsilon)$\hfill .\,.\,.} & $9{,}686\,2484$\\[.3em]
\makebox[18em][l]{$\log\cos\tfrac{1}{2}(i+\epsilon)$\hfill .\,.\,.} & $9{,}941\,6108$
\end{tabular}
\end{center}
\begin{center}
\begin{tabular}{@{}ll@{}}
$\tfrac{1}{2}(\AscendingNode-n)=$ & $-3^0\,45^{\prime}\,53^{\prime\prime}{,}15$\\[.3em]
$\tfrac{1}{2}(i+\epsilon)=$ & $29\quad2\quad58{,}45$\\[.3em]
$\tfrac{1}{2}(i-\epsilon)=$ & $5\quad35\quad2{,}65$\\[.3em]
\makebox[18em][l]{$\log\cos\tfrac{1}{2}(\AscendingNode-n)$\hfill .\,.\,.} & $9{,}999\,0618$\\[.3em]
\makebox[18em][l]{$\log\sin\tfrac{1}{2}(i-\epsilon)$\hfill .\,.\,.\,.} & $8{,}988\,1405$\\[.3em]
\makebox[18em][l]{$\log\cos\tfrac{1}{2}(i-\epsilon)$\hfill .\,.\,.\,.} & $9{,}997\,9342$
\end{tabular}
\end{center}
\Needspace{7\baselineskip}
\noindent Daraus folgt
\begin{center}
\begin{tabular}{@{}ll@{}}
\makebox[18em][l]{$\log\sin\tfrac{1}{2}i'\sin\tfrac{1}{2}(\AscendingNode'+\Delta)$\hfill .\,.\,.} & $8{,}503\,5510_n$\\[.3em]
\makebox[18em][l]{$\log\sin\tfrac{1}{2}i'\cos\tfrac{1}{2}(\AscendingNode'+\Delta)$\hfill .\,.\,.} & $8{,}987\,2023$\\[.3em]
\hline
daraus $\tfrac{1}{2}(\AscendingNode'+\Delta)=$ & $341^0\,49^{\prime}\,19^{\prime\prime}{,}01$\\[.3em]
\makebox[18em][l]{$\log\sin\tfrac{1}{2}i'$\hfill .\,.\,.\,.\,.\,.\,.\,.\,.\,.\,.\,.\,.} & $9{,}009\,4368$
\end{tabular}
\end{center}
\begin{center}
\begin{tabular}{@{}ll@{}}
\makebox[18em][l]{$\log\cos\tfrac{1}{2}i'\sin\tfrac{1}{2}(\AscendingNode'-\Delta)$\hfill .\,.\,.} & $8{,}758\,9134_n$\\[.3em]
\makebox[18em][l]{$\log\cos\tfrac{1}{2}i'\cos\tfrac{1}{2}(\AscendingNode'-\Delta)$\hfill .\,.\,.} & $9{,}996\,9960$\\[.3em]
\hline
daraus $\tfrac{1}{2}(\AscendingNode'-\Delta)=$ & $356^0\,41^{\prime}\,31^{\prime\prime}{,}43$\\[.3em]
\makebox[18em][l]{$\log\cos\tfrac{1}{2}i'$\hfill .\,.\,.\,.\,.\,.\,.\,.\,.\,.\,.\,.\,.} & $9{,}997\,7202$
\end{tabular}
\end{center}
\noindent Wir erhalten somit $\frac{1}{2}i'=5^0\,51^{\prime}\,56^{\prime\prime}{,}445$, $i'=11^0\,43^{\prime}\,52^{\prime\prime}{,}89$,\hfil\break $\AscendingNode'=338^0\,30^{\prime}\,50^{\prime\prime}{,}44$, $\Delta=-14^0\,52^{\prime}\,12^{\prime\prime}{,}42$. Der Punkt $n$ auf der Himmelskugel entspricht übrigens offenbar- dem Herbstpunkt. Somit beträgt der Abstand des aufsteigenden Knotens der Bahn auf dem Äquator vom Frühlingspunkt (seine \textit{Rektaszension}) $158^0\,30^{\prime}\,50^{\prime\prime}{,}44$.

\SourceLeaf{075}Zur Veranschaulichung von Art. 53 führen wir dieses Beispiel weiter und entwickeln Formeln für die Koordinaten in Bezug auf drei Ebenen durch die Sonne. Eine davon sei zum Äquator parallel, und die positiven Pole der beiden übrigen Ebenen mögen bei den Rektaszensionen $0^0$ und $90^0$ liegen. Die Abstände von diesen Ebenen seien der Reihe nach $z$, $x$ und $y$. Bezeichnen wir außerdem den Abstand des heliozentrischen Orts auf der Himmelskugel von den Punkten $\AscendingNode$ und $\AscendingNode'$ jeweils mit $u$ und $u'$, so gilt $u'=u-\Delta=u+14^0\,52^{\prime}\,12^{\prime\prime}{,}42$. Den in Art. 53 mit $i$, $N-\AscendingNode$ und $u$ bezeichneten Größen entsprechen hier $i'$, $180^0-\AscendingNode'$ und $u'$. Aus den dort angegebenen Formeln folgt damit
\begin{center}
\begin{tabular}{@{}ll@{}}
\makebox[18em][l]{$\log a\sin A$\hfill .\,.\,.\,.} & $9{,}968\,7197_n$\\[.3em]
\makebox[18em][l]{$\log a\cos A$\hfill .\,.\,.\,.} & $9{,}554\,6380_n$\\[.3em]
\hline
daraus $A=$ & $248^0\,55^{\prime}\,22^{\prime\prime}{,}97$\\[.3em]
\makebox[18em][l]{$\log a$\hfill .\,.\,.\,.\,.\,.} & $9{,}998\,7923$
\end{tabular}
\end{center}
\begin{center}
\begin{tabular}{@{}ll@{}}
\makebox[18em][l]{$\log b\sin B$\hfill .\,.\,.\,.} & $9{,}563\,8058$\\[.3em]
\makebox[18em][l]{$\log b\cos B$\hfill .\,.\,.\,.} & $9{,}959\,5519_n$\\[.3em]
\hline
daraus $B=$ & $158^0\,5^{\prime}\,54^{\prime\prime}{,}97$\\[.3em]
\makebox[18em][l]{$\log b$\hfill .\,.\,.\,.\,.\,.} & $9{,}992\,0848$
\end{tabular}
\end{center}
\noindent Wir erhalten somit
\[
\begin{aligned}
x&=ar\sin(u'+248^0\,55^{\prime}\,22^{\prime\prime}{,}97)\\
&=ar\sin(u+263^0\,47^{\prime}\,35^{\prime\prime}{,}39)\\[.4em]
y&=br\sin(u'+158\quad5\quad54{,}97)\\
&=br\sin(u+172\quad58\quad7{,}39)\\[.4em]
z&=cr\sin u'\\
&=cr\sin(u+14\quad52\quad12{,}42),
\end{aligned}
\]
wobei $\log c=\log\sin i'=9{,}308\,1870$ ist.

Eine andere Lösung der hier behandelten Aufgabe findet sich bei \textsc{von Zach} in der \textit{Monatliche Correspondenz}, Bd. IX, S. 385\HistoricalNoteEight{Mai 1804. Band VI. S. 94.]}.
